\documentclass[11pt,oneside]{article}
\pdfoutput=1

% ============================================================
%  The Conceptual Space of Software Architecture
%  The Latent Variables That Generate Architectural Decisions
% ============================================================

\usepackage[utf8]{inputenc}
\usepackage[T1]{fontenc}
\usepackage{lmodern}
\usepackage[a4paper,margin=1in]{geometry}
\usepackage{amsmath,amssymb,amsthm}
\usepackage{microtype}
\usepackage{enumitem}
\usepackage{booktabs}
\usepackage{longtable}
\usepackage{array}
\usepackage{xcolor}
\usepackage{titlesec}
\usepackage{fancyhdr}
\usepackage{setspace}
\usepackage[hidelinks,breaklinks]{hyperref}
\usepackage{parskip}

\definecolor{accent}{RGB}{40,72,120}
\hypersetup{colorlinks=true,linkcolor=accent,citecolor=accent,urlcolor=accent}

\titleformat{\section}{\Large\bfseries\color{accent}}{\thesection}{0.6em}{}
\titleformat{\subsection}{\large\bfseries}{\thesubsection}{0.6em}{}
\titleformat{\subsubsection}{\normalsize\bfseries\itshape}{\thesubsubsection}{0.6em}{}

\pagestyle{fancy}
\fancyhf{}
\fancyhead[L]{\small\itshape The Conceptual Space of Software Architecture}
\fancyhead[R]{\small\thepage}
\renewcommand{\headrulewidth}{0.4pt}

\newtheorem{claim}{Claim}
\newtheorem{prop}{Proposition}
\newtheorem{defn}{Definition}
\theoremstyle{definition}
\newtheorem{axis}{Axis}
\newtheorem{pred}{Prediction}

\setstretch{1.06}

\title{\vspace{-1.0cm}\textbf{\color{accent}The Conceptual Space of Software Architecture}\\[0.35cm]
\large The Latent Variables That Generate Architectural Decisions}
\author{Edwin O Onyango\\[-1pt]
\small Dartmouth College\\
\small \texttt{edwin.o.onyango.jr@gmail.com}
\and
Walter Wagunde\\[-1pt]
\small Tufts University}
\date{August 2026}

\begin{document}
\maketitle
\thispagestyle{empty}

\begin{abstract}
\noindent What is software architecture, such that a machine could hold it? Every coding agent
today navigates a codebase through its observables: the file tree, a symbol graph, text
similarity, and a frozen prose file in which humans write down what actually matters. This paper
asks whether the thing being written down has a small generating basis. By that we mean a short
list of latent variables such that every architectural decision, in any system, is a move along
one of them, the way physics replaced a bestiary of phenomena with a few forces.

We answer in three ways. The first is theoretical. We start from one premise: software is kept
alive by a bounded population of agents, human or machine, while the world changes around it.
From that premise we derive what an architecture must be. It is not a structure. It is a set of
decisions laid over four graphs. The problem has a graph of constraints. The code has a graph
of dependencies. The runtime has a graph of interactions. The organization has a graph of
communication. We then read sixty years of theory, from Parnas, Dijkstra, Simon, and Conway
through design structure matrices, architectural styles, and architecture as decisions, and we
find the same small set of variables appearing again and again under different names. Where do
the cuts fall on each graph? Do the four cuts line up with one another? What kind of
interaction crosses each boundary? Which orders of visibility and abstraction sit over the
cuts? And for each decision: how volatile is it, when does it take force, what would reversing
it cost, and how far does it reach? One such agreement between two theories could be a
coincidence. What we find is a pattern of them. The measure now called propagation cost was
derived twice, once from a change-cost argument and once from a build-cost argument, by authors
who did not cite each other. Binding time appears independently in five separate frameworks.
The dominance order was discovered four times. When many lines of work that start from
different premises and share no vocabulary keep arriving at the same variables, the most
reasonable explanation is that the variables are really there. That logic of convergent
evidence is the backbone of the paper.

The second way is empirical. We survey eighteen major systems and five industrial migrations,
and of every architectural difference we find, we ask one simple question: is this a choice,
or is it decay? Some differences are choices. Excellent teams pick opposite answers to the
same question, each side for good reasons, and over the years real systems move in both
directions. Other differences are not choices at all. Nobody chooses a tangle of circular
dependencies; systems slide into it, and climbing out costs effort. The paper keeps these two
kinds strictly apart, calling the choices coordinates and the decay directions potentials.
The clearest single lesson from the migrations is also simple. How you divide your code and
how you deploy it are two separate decisions, and the teams that treated them as one are the
ones that later wrote post mortems about it; three of five name exactly this as the root
error. The third way is computational. Roughly three quarters of the variables can be
estimated today from a bare repository or its git history. Only the intent-carrying ones
require a human, and for those the reflexion model loop is a proven protocol: the machine
drafts, the human corrects, the machine enforces.

Finally, we assemble the basis into a conceptual space in G\"ardenfors' sense, the picture
cognitive science uses for how minds hold concepts: qualities become dimensions, concepts
become regions, similarity becomes nearness. Styles are its prototype regions. Tactics and the
six modular operators are its move set. Quality attributes form a value field, held separate
from the coordinates. The metric is asymmetric and built on reachability. That choice turns
Tversky's classical objection, that human judgments of similarity refuse to behave like
distances, from a problem into a prediction.
We are equally explicit about what the space is not. It is not fixed, since framing re-weights it. It is not Euclidean, since edge
costs are nonlocal. It is not one partition, since recovery algorithms disagree precisely because
each optimizes a different real dimension. We close with falsifiable predictions and with a
construction recipe: what an agent must compute so that, on seeing files, functions, and classes,
it is aware of every variable in play.
\end{abstract}

\tableofcontents
\vspace{0.4cm}
\hrule

% ===========================================================================
\section{The question, sharpened}
% ===========================================================================

We want to begin somewhere that almost feels too obvious to mention, because it is usually in
those places that the most important assumptions sit unnoticed.

Ask two engineers to draw the architecture of the system they share and you will get two
different diagrams, drawn at different granularities, organized around different concerns. Now
ask those same two engineers which proposed change is dangerous, which module is load-bearing,
and which shortcut will be regretted, and they will largely agree. Whatever the architecture is,
it is not the diagrams. It is the thing that lets two people who have never compared diagrams
make the same predictions. This paper is an attempt to pin that thing down, precisely enough
that a machine could estimate it.

There is a weak version of this research and a strong one, and only the strong one is worth
doing.

The weak version says that architectures differ in many ways, and that we should catalog the
ways: styles, patterns, tactics, smells. This has been done, repeatedly and well
\cite{shawgarlan1996,gamma1994,bass2012}, and it is nearly useless for our purpose, because a
catalog has no generative power. It tells you what has been seen, not what must be. A catalog
cannot tell an agent, facing a system it has never seen, which questions to ask of it.

The strong version is a claim about generating variables. By latent we mean nothing exotic, only
the statistician's plain sense of the word: quantities never observed directly, whose existence
is earned by how much of the observable record they explain. The claim is that there exists a
small basis $B$ of such quantities, and that every architectural decision, from ``should this be
two services or one'' down to ``should this function know about that type,'' is a move along the
elements of $B$. Every architecture is then a point positioned by them, or more precisely, as we
will see, a trajectory. If $B$ exists and is small, then understanding a codebase is not
memorizing its shape. It is estimating its coordinates. ``Know one system, understand many''
stops being an aspiration and becomes a structural consequence, because an observation about one
boundary updates a shared latent and moves expectations about every other boundary.

This is the same move made twice before in adjacent territory. Physics replaced phenomena with
forces. Machine learning replaced hand-written rules with latent representations that earn their
existence by making prediction easier. The object here is different, and choosing it correctly
is most of the work. We are not asking what architecture documentation contains; that question
produces templates. We are not asking what architecture recovery outputs; that question produces
partitions. We are asking: \textbf{what are the degrees of freedom of the design, such that
fixing their values determines every structural fact about a system that an engineer would call
architectural, and such that every architectural disagreement is a disagreement about their
values?}

The question has an unusual property. It has been answered many times, partially, by people who
mostly did not read each other. Parnas answered it with the distribution of anticipated change
\cite{parnas1972}. Dijkstra answered it with ordered abstraction \cite{dijkstra1968}. Simon
answered it with interaction density \cite{simon1962}. Conway answered it with the shape of the
designing organization \cite{conway1968}. Baldwin and Clark answered it with option value over a
dominance hierarchy \cite{baldwinclark2000}. Shaw and Clements answered it with an explicit
dimension table \cite{boxology1997}. Fielding answered it by deriving an architecture as a walk
through constraint space \cite{fielding2000}. Jansen and Bosch answered it by declaring the
decisions themselves to be the architecture \cite{jansenbosch2005}. The scandal of the field is
not that no theory exists. It is that eight partial theories exist, and nobody has computed
their union, checked it against the decision points of real systems, and asked what is
irreducible.

That is this paper's job. Its thesis, stated now and defended throughout:

\begin{claim}\label{claim:main}
Software architecture is a decision structure laid over four graphs. Its latent variables fall
into four tiers: (i) \textbf{boundary geometry}, which says where the cuts fall on each graph,
how the four cuts align, and what kind of interaction crosses each boundary; (ii) \textbf{order
structure}, the partial orders of use, abstraction, and dominance laid over the cuts; (iii)
\textbf{a commitment calculus}, which attaches to each decision its volatility, binding time,
reversal cost, blast radius, and commitment state; and (iv) \textbf{a value field and a
dynamics}, quality-attribute gradients held separate from the coordinates, and an entropic drift
that makes architecture a trajectory rather than a point. These tiers are mutually irreducible,
jointly generative, and largely machine-estimable.
\end{claim}

The rest of the paper earns each clause of that sentence.

% ===========================================================================
\section{Method: three triangulating tests}
% ===========================================================================

A proposed latent variable is cheap. The literature contains hundreds of candidate ``dimensions
of architecture,'' and admitting them by assertion would only add one more catalog to the pile.
We admit a variable into the basis only when it passes tests of a different kind. Three tests do
the work throughout this paper.

\subsection{The convergence test}

If a quantity is a real dimension of the space rather than an artifact of one author's
formalism, it should be discovered independently, from different premises, by people solving
different problems. The flagship example is the density of the transitive closure of the
dependency graph. MacCormack, Rusnak and Baldwin derived it as propagation cost, the expected
fraction of a system affected by a random change \cite{maccormack2006}. A decade earlier, Lakos
had derived it as average cumulative component dependency, the expected cost of an incremental
build and test \cite{lakos1996}. Same mathematical object, reached once from a change-cost
argument and once from a build-cost argument. When two theories that share no citations and no
vocabulary force the same coordinate, the coordinate is probably in the territory, not the map.
And one such convergence could still be a coincidence, which is why the method is cumulative.
We collect convergences until coincidence stops being a plausible explanation. Binding time,
for example, appears independently in five separate frameworks (Axis~\ref{ax:binding}), and
the dominance order was discovered four times (Axis~\ref{ax:dominance}). We will flag every
such convergence as it arrives. There are more of them than the field seems to know.

\subsection{The two-way traffic test}

Faced with any difference between two systems, we ask one simple question. Is this a choice, or
is it decay? A choice is a genuine degree of freedom. Competent designers may sit at either
pole, depending on the forces they face, and we will call such an axis a dimension, or a
coordinate. Decay is different. It is a direction no competent designer takes on purpose, a
slope that systems slide down and must spend effort to climb back up, and we will call it a
potential. A catalog cannot tell these two apart. The empirical record can. If respected
systems occupy opposite poles of an axis with explicit, opposite justifications, and stronger
still, if documented migrations exist in both directions, then the axis is a coordinate. Chromium moved its isolation boundary finer after Spectre
\cite{chromium2009,reis2019}. Istio merged its control plane coarser after measuring operational
cost \cite{istio2020}. Segment went from monolith to microservices and back to monolith
\cite{segment2018}. Apache and nginx converged on each other's concurrency mechanisms from
opposite ends \cite{aosanginx}. Isolation granularity, centralization, binding time, and
concurrency model all carry two-way traffic. By contrast, nobody migrates toward a larger cyclic
core of mutually dependent files on purpose. Dependency-graph health carries one-way intentional
traffic only, and it therefore enters the theory not as a coordinate but as a potential, a
gradient that the dynamics descend or resist (\S\ref{sec:dynamics}). This distinction between
coordinates and potentials, which the two-way test operationalizes, is one of the paper's
load-bearing contributions.

\subsection{The computability test}

The end state of this research program is an agent that, on opening a repository, estimates the
variables. A latent variable that is well defined but unestimable in principle, whether from
code, history, traces, or a bounded human interaction, can still constrain theory, but it cannot
enter the working basis. We therefore tag every variable with its evidence requirement:
\emph{bare repo}, \emph{plus git history}, \emph{plus runtime traces}, or \emph{plus human
intent}. And we take seriously the reflexion-model result that the human-intent inputs are
bounded and cheap. One engineer reflexion-modeled the 1.2 MLOC Excel codebase in about a month,
iterating a hypothesized module map against computed divergences \cite{murphy1995}. The machine
drafts, the human corrects, the machine enforces. Intent is an input, not an excuse.

\subsection{What we deliberately keep separate}

Three separations are enforced everywhere, because their conflation is the commonest failure of
prior syntheses. The first is coordinates versus value field. Where a system \emph{is}, its
structural position, is one thing. How \emph{good} that position is for given goals is another.
Lane's design spaces \cite{lane1990}, ATAM's sensitivity points \cite{atam2000}, and Fielding's
property table \cite{fielding2000} are all maps between these two objects, and they lose their
content if the objects are merged. The second is coordinates versus dynamics. The drift of
structure under change, Lehman's laws \cite{lehman1980}, is a vector field over the space, not
an axis of it. The third is the four graphs themselves. Requirements, code, runtime, and
organization each carry their own cut, and much of what the field argues about is secretly a
claim about the congruence of two of them (\S\ref{sec:fourgraphs}).

% ===========================================================================
\section{The object: what an architecture is}\label{sec:object}
% ===========================================================================

We now derive the object whose space we are charting, from premises weak enough to be
undeniable.

\textbf{Premise.} A software system is an artifact maintained by a bounded population of agents,
human or machine, under sustained and partially unanticipated change, for purposes located
outside the artifact.

Everything follows from the three pressures in that sentence: \emph{boundedness}, since no agent
holds the whole; \emph{change}, since the artifact is never done, which is Lehman's E-type
condition \cite{lehman1980}; and \emph{external purpose}, since validity is judged in a moving
world rather than against a frozen specification.

\subsection{Boundedness forces cuts; change decides where they go}

Because no agent holds the whole, the system must be cut into parts small enough to hold. This
is Simon's argument that complex systems built by bounded processes are hierarchies of nearly
decomposable parts, with interaction density high within parts and low between them
\cite{simon1962,simonando1961}. Throughout this paper a \emph{cut} is nothing exotic. It is a
partition of a graph's vertices into blocks, a decision about where the lines get drawn.

But density of \emph{which} interactions? Here sits the first deep fork in the classical
literature, and it is worth slowing down for, because it is not a disagreement to be resolved.
It is a pair of coordinates to be kept distinct. Structured design cuts the graph of data-flow
connections \cite{stevens1974}. Alexander cuts the graph of requirement interactions
\cite{alexander1964}. Simon cuts the graph of runtime interactions. Parnas rejected all
flowchart-shaped criteria and cut by anticipated change: begin ``with a list of difficult design
decisions or design decisions which are likely to change,'' and make each one the hidden secret
of one module \cite{parnas1972}. Interaction minimization and change localization coincide only
under an assumption that nobody in the classical corpus states or tests: that change propagates
along interaction edges with uniform probability. The KWIC construction is precisely a system
where they disagree. Parnas's Decomposition 1 scores tolerably on coupling and fails
catastrophically on change containment. \emph{The volatility field over decisions is therefore
an independent input, not derivable from any interaction graph} (\S\ref{sec:commitment}).

\subsection{The four graphs}\label{sec:fourgraphs}

Collecting the graphs that the classical theories partition, we get exactly four, and the list
appears to be complete:

\begin{enumerate}[leftmargin=1.6em]
\item $G_P$, the \textbf{problem graph}: requirement and misfit variables and their interactions
(Alexander's constraint graph \cite{alexander1964}; feature models; the domain's own conceptual
joints, which domain-driven design later named bounded contexts \cite{evans2003}).
\item $G_C$, the \textbf{code graph}: source units and their syntactic dependencies (the DSM
substrate \cite{steward1981,maccormack2006}; the uses relation \cite{parnas1979}).
\item $G_R$, the \textbf{runtime graph}: executing elements and their interactions (Fielding's
components, connectors, and data \cite{fielding2000}; process and deployment topology).
\item $G_O$, the \textbf{organization graph}: the agents and their communication structure
(Conway \cite{conway1968}).
\end{enumerate}

Each classical theory is, at bottom, the same mathematical move applied to a different member of
this family: partition a weighted graph so that heavy edges fall inside blocks. The moment the
family is seen whole, a new class of variables appears that no single-graph theory can express.
These are the \textbf{congruence terms}, the degree of alignment between cuts on different
graphs. At least three are independently load-bearing:

\begin{itemize}[leftmargin=1.6em]
\item $\kappa_{OC}$, the congruence of organization and code, which is Conway alignment. This is
the best-evidenced variable in the entire literature. Organizational metrics predicted
failure-prone Vista binaries better than any code metric \cite{nagappan2008}. In five matched
pairs of functionally equivalent products, the more distributed organization produced the more
modular codebase in every pair \cite{maccormack2012}. The concentration of ownership at a
boundary predicts its failure rate \cite{bird2011}. And loosely coupled architecture,
operationalized as teams able to test and deploy without coordination, is among the strongest
population-scale predictors of delivery performance \cite{accelerate2018}. The mirroring is a
prevalent pattern, not a law. It weakens in open source \cite{colferbaldwin2016}, which is
exactly what makes it a \emph{variable}.
\item $\kappa_{PC}$, the congruence of problem and code: whether module boundaries track the
problem's own constraint clusters. This was Alexander's original program \cite{alexander1964},
and DDD's bounded contexts are its modern operationalization \cite{evans2003}.
\item $\kappa_{CR}$, the congruence of the logical and the physical: whether code-level module
boundaries and runtime deployment boundaries coincide. The empirical record of
\S\ref{sec:empirical} shows this is the single most expensive conflation in industrial practice.
Modularity forces, meaning ownership, comprehension, and change isolation, act on the logical
cut. Operational forces, meaning cost, latency, consistency, and blast radius, act on the
physical cut. Three of five major migration post-mortems name their root error as having coupled
the two \cite{primevideo2023,shopify2019,uberdoma2020}.
\end{itemize}

\subsection{Decisions, not structures, are the substance}

A partition of four graphs is still a static description. Two further classical results force
the final reconceptualization. Parnas's program-families paper models development as a tree
whose nodes are incomplete designs and whose edges are design decisions. An architecture is a
node, a prefix-closed set of commitments, and a family is a subtree \cite{parnas1976}. Thirty
years later, and apparently without genealogical connection, Jansen and Bosch re-derived the
same identification from maintenance economics. Structure is the residue. The decisions, with
their rules and constraints, are the architecture, and their loss, which they called knowledge
vaporization, is the causal mechanism of design erosion \cite{jansenbosch2005,vangurp2002}.
Taylor, Medvidovic and Dashofy made it the textbook definition: architecture is the set of
principal design decisions \cite{taylor2010}. The convergence test passes. \emph{The point of
the space is a decision set, and the structural facts, the cuts and the orders, are its shadow
on the four graphs.}

This identification does real work. It explains why architecture cannot be fully recovered from
code: bans and rejected alternatives leave no residue, which is why Kruchten's ontology gives
ban decisions, crucial and invisible, a first-class place \cite{kruchten2004}. It explains why
erosion is path dependence plus amnesia rather than bad decisions \cite{vangurp2002}. And it
explains why the space's metric must be asymmetric (\S\ref{sec:metric}): moving between two
architectures is not a distance between shapes but a re-decision cost, and re-deciding is
directional.

\begin{defn}[Architecture]\label{def:arch}
An architecture is a pair $(D, \sigma)$ where $D$ is a set of design decisions carrying the
commitment calculus of \S\ref{sec:commitment} and partially ordered by their restriction
relationships, and $\sigma$ is the induced structure: the cuts $\Pi_P, \Pi_C, \Pi_R, \Pi_O$ on
the four graphs, the congruences $\kappa$ among them, the partial orders of \S\ref{sec:orders}
over the cuts, and the interaction type at each boundary (Axis~\ref{sec:boundarytype}). A
codebase realizes an architecture with a conformance residual (reflexion's convergences,
divergences, and absences \cite{murphy1995}; erosion and drift are the residual's named
pathologies \cite{perrywolf1992}).
\end{defn}

% ===========================================================================
\section{The basis, tier by tier}\label{sec:basis}
% ===========================================================================

We now assemble the basis. Each variable is stated as an axis: what varies, its formal
structure, the independent discoveries that force it, which is the convergence test at work, and
its evidence requirement, which is the computability test at work. The four tiers of
Claim~\ref{claim:main} organize the presentation.

\subsection{Tier I: boundary geometry}

The first tier answers the plainest questions about a system. Where are the lines drawn, on
which graph, and at what scale? Do the lines on different graphs agree? And what actually
crosses them?

\begin{axis}[Cut placement, per graph]\label{ax:cut}
For each graph $G_i \in \{G_P, G_C, G_R, G_O\}$: the partition $\Pi_i$ of its vertices, with
Simon's near-decomposability ratio $\varepsilon_i$, the ratio of inter-block to intra-block
interaction density, as its principal scalar projection. \emph{Formal structure:} the lattice of
partitions; nearly block-diagonal matrices, with a proved timescale-separation theorem for the
linear case \cite{simonando1961}. \emph{Convergences:} Simon on runtime interactions; Alexander
on requirement interactions, with an explicit min-cut algorithm \cite{alexander1964}; structured
design on data-flow connections \cite{stevens1974}; community detection and every
software-clustering objective since \cite{mancoridis1999}. \emph{Evidence:} bare repo for
$G_C$; traces for $G_R$; authorship and communication records for $G_O$; requirements analysis
or domain modeling for $G_P$.
\end{axis}

\begin{axis}[Granularity, per graph]\label{ax:granularity}
The characteristic scale of the cut: statement to function to class to module to service to
system on $G_C$; thread to process to sandboxed process to container to node on $G_R$;
individual to team to department on $G_O$. Categorization research predicts, and practitioner
discourse is consistent with, a paradigm-relative \emph{basic level} at which practitioners
reason by default: the module in monolith cultures, the service in distributed ones. Granularity
is therefore not only a structural choice. It sets the resolution of the community's default
vocabulary \cite{rosch1978,gardenfors2000}. \emph{Two-way traffic:} Chromium moved isolation
finer, first process-per-site-instance \cite{chromium2009}, then strict site isolation after
Spectre changed the threat model \cite{reis2019}; Istio consolidated four control-plane
deployables into one \cite{istio2020}. Both were competent; the forces differed.
\emph{Evidence:} bare repo and deployment manifests.
\end{axis}

\begin{axis}[Congruence terms]\label{ax:congruence}
$\kappa_{OC}, \kappa_{PC}, \kappa_{CR}$ as introduced in \S\ref{sec:fourgraphs}: matrix
congruence statistics between cuts on different graphs. These are second-order variables,
properties of \emph{pairs} of cuts, and they are where several of the field's most robust
empirical results actually live \cite{nagappan2008,maccormack2012,cataldo2008,accelerate2018}.
\emph{Evidence:} $\kappa_{OC}$ from git authorship crossed with dependency structure;
$\kappa_{CR}$ from build and deploy manifests crossed with module structure.
\end{axis}

\begin{axis}[Boundary interaction type]\label{sec:boundarytype}
At every cut edge, the interaction is not a scalar ``coupling strength.'' It is a structured
product of six factors, enumerated independently by the boxology program \cite{boxology1997}
and the connector taxonomy \cite{mehta2000}. The first factor is \emph{control topology}
crossed with \emph{data topology}: each may be linear, acyclic, hierarchical, star, or
arbitrary, they are stated separately, and a relation records whether they coincide and which
way they point. The second is \emph{synchronicity}: lockstep, synchronous, asynchronous, or
opportunistic. The third is \emph{data mode}: passed, shared, copy-out-copy-in, or broadcast.
The remaining three are \emph{delivery guarantees}, \emph{cardinality}, and \emph{state locus},
meaning where conversation state lives: in a server session, held by the client, distributed,
or nowhere \cite{fielding2000}.\footnote{A sourcing caveat applies to the enumerated value sets of the
boxology, Lane, connector-taxonomy, and tactics tables throughout this paper: the primary
technical reports sit behind access restrictions, and the value sets are reconstructed from the
published texts and secondary quotations. The dimensional \emph{structure} is not in doubt;
individual value lists remain to be verified against the primary tables.} The classical coupling
ladder (content, common, external, control, stamp, data, in increasing order of virtue,
assembled from ingredients in \cite{stevens1974} and codified by the structured-design school
\cite{yourdon1979}) is an ordinal compression of this product: historically first, dimensionally
coarse. \emph{Evidence:} bare repo for interface signatures and connector idioms; traces for the
realized subset.
\end{axis}

\begin{axis}[Interface thinness and uniformity]\label{ax:interface}
Per boundary: how much of the hidden side is deducible through it. This is Parnas's ``reveal as
little as possible'' \cite{parnas1972}, with three independently derived formalizations. Simon
argued that only slow, aggregate quantities should cross a near-decomposable boundary
\cite{simon1962}. Baldwin's thin crossing points locate transactions where transfers are few and
standardizable \cite{baldwin2008}. The information-theoretic reading counts the bits of the
secret leaked by the interface \cite{allen1999}. To thinness add \emph{uniformity}: bespoke
versus domain-standard versus universal, as in Unix pipes and HTTP verbs, which Fielding shows
trades efficiency for evolvability and visibility \cite{fielding2000}. And add \emph{stability
asymmetry}: a system may expose stable narrow seams over volatile interiors, as in Git's
plumbing and porcelain, or Linux's stable ABI outward and unstable API inward, making interface
granularity independent of component granularity \cite{aosagit}. \emph{Evidence:} bare repo.
\end{axis}

\subsection{Tier II: order structure}\label{sec:orders}

Cuts say where the lines are. The second tier says which way is up: who may know about whom, who
sits above whom in interpretation, and whose decisions bind whose.

The classical corpus contains a standing warning here, first issued in Parnas's 1972 paper
itself: at least three distinct partial orders masquerade as ``hierarchy,'' and collapsing them
is a category error \cite{parnas1972}.

\begin{axis}[The uses order]\label{ax:uses}
The designed relation ``$A$ uses $B$,'' meaning the correctness of $A$ requires a correct $B$,
which is not the same as invokes. Its shape ranges from unconstrained digraph to strict layers.
Its value is exact by construction: the downward-closed sets of the uses order are precisely the
buildable, deliverable subsets of the system, so extension and contraction capability is the
lattice of order ideals \cite{parnas1979}. Simon's stable intermediate forms, the watchmaker
argument \cite{simon1962}, are the same object seen dynamically. \emph{Evidence:} bare repo.
\end{axis}

\begin{axis}[The abstraction order]\label{ax:abstraction}
Dijkstra's levels: each level implements an abstract machine above which certain resources
``lose their identity.'' Formally, an ordered chain of state-space quotients, whose payoff is
combinatorial, since verification cost falls from the product of level state-spaces to their sum
\cite{dijkstra1968}. This is distinct from the uses order. A uses-hierarchy need not be a
hierarchy of interpretation, and the two orders can disagree on the same system.
\emph{Evidence:} bare repo, partially; abstraction boundaries are inferable from type structure
and naming, confirmable by a human.
\end{axis}

\begin{axis}[The dominance order and design rules]\label{ax:dominance}
Baldwin and Clark's deepest structural claim: a modularization is not a partition alone but a
partition \emph{plus a visibility hierarchy}. Design rules are decisions fixed early and visible
to all, and hidden parameters may not constrain them \cite{baldwinclark2000}. Cai and Sullivan
showed the DSM is not primitive but the shadow of a constraint network plus this dominance order
\cite{cai2006}. Its measurable projections are the design rule hierarchy and the decoupling
level metric, evaluated on 129 projects with maintainability validation on a 41-project subset
\cite{mo2016}, though the validations to date come largely from one research group, a
replication gap the theory should not paper over. Martin's stability direction, with instability
decreasing along dependency edges, and the refinement orders of interface and contract theories
(alternating simulation \cite{dealfaro2001}; assume-guarantee contracts \cite{benveniste2018})
are the same order structure in other clothing. Four convergent discoveries. \emph{Evidence:}
bare repo.
\end{axis}

\begin{axis}[Dependency orientation relative to the abstraction gradient]\label{ax:orientation}
Whether dependencies flow toward abstraction and stability or against it. This is the
main-sequence discipline, SDP and SAP \cite{martin2002}, with hexagonal, onion, and clean
architecture as its packaging \cite{cockburn2005}. On Martin's $(I,A)$ unit square, instability
against abstractness, the discipline is literally a vector-field constraint with a named
prototype region, the main sequence $A + I = 1$, and named pathological corners. It is the most
literal small conceptual space in the existing literature, and a template for this paper's
construction. \emph{Evidence:} bare repo, for typed languages.
\end{axis}

\begin{axis}[Reachability density]\label{ax:reach}
The scalar that summarizes what the orders jointly control: the density of the transitive
closure of $G_C$. This is propagation cost \cite{maccormack2006}, and it is Lakos's ACD/N
\cite{lakos1996}, the paper's flagship convergence. Its refinement, decoupling level, discounts
legitimate dependence on design rules and can rank systems oppositely to raw propagation cost.
The two are different projections, and the space needs both \cite{mo2016}. Derived
classifications give the space some of its best-attested prototype regions: core size,
core-periphery versus hierarchical versus multi-core types, element roles on the (VFI, VFO)
plane. Most large real systems are core-periphery \cite{baldwin2014}. One caveat matters for the
geometry: the cost of an edge is nonlocal, since a single added edge can double the transitive
closure, so the space is non-Euclidean in the edge basis \cite{lakos1996}. \emph{Evidence:}
bare repo.
\end{axis}

\subsection{Tier III: the commitment calculus}\label{sec:commitment}

Tiers I and II describe the shadow. Tier III is the substance: the variables attached to the
decisions themselves. The decision-centric literature supplies them, from Parnas 1976 through
Jansen and Bosch, Kruchten, the rationale tradition, and the options tradition, and they compose
into what we call the commitment calculus.

\begin{axis}[The volatility field]\label{ax:volatility}
An assignment of anticipated-change weight to every design decision. This is Parnas's generative
criterion \cite{parnas1972}, factored at scale into change-source classes of hardware, behavior,
and internal \cite{parnas1985}, and re-derived economically as per-module technical potential
$\sigma_i$ in the options theory, where volatility is where option value lives
\cite{baldwinclark2000,sullivan2001}. The two sign conventions, volatility as hazard in the
stability principles and volatility as resource in the options tradition, agree on the
structural prescription that nothing should depend on the volatile, and disagree on valuation.
The space needs the position and the sign kept separate. \emph{Evidence:} git history, churn and
co-change, as proxy; requirements volatility needs human input.
\end{axis}

\begin{axis}[Binding time]\label{ax:binding}
When each decision takes force: design, compile, link, startup, or run time. This appears
independently in boxology, stated separately for control and data partners \cite{boxology1997};
in the connector taxonomy, as linkage binding \cite{mehta2000}; in the modifiability tactics, as
defer binding \cite{bass2012}; in program families, where family-common decisions bind early and
member-variable decisions bind late \cite{parnas1976}; and across the empirical corpus, where
nginx's compile-time modules and Eclipse's runtime OSGi bundles were chosen oppositely under
explicit forces \cite{aosanginx,aosaeclipse}. Five independent appearances. This is the
strongest convergence in the paper after reachability density. \emph{Evidence:} bare repo.
\end{axis}

\begin{axis}[Reversal cost, blast radius, and accrued dependents]\label{ax:reversal}
Fowler defined architecture, with a deliberate wink, as the things people perceive as hard to
change, while insisting the architect's job is to shrink that set \cite{fowler2003}. The
definition decomposes into a product of terms. The first is reversal cost: what a rewrite, a
migration, or data gravity would demand. The second is blast radius: Kruchten's scope of
system, temporal, and organizational extent \cite{kruchten2004}. The third is dependents
accrued, the in-degree in the decision graph, because irreversibility often arrives after the
decision, by accretion. Binding time modulates all three. The options literature
prices deferral against exactly these \cite{baldwinclark2000}. The one-way and two-way door
rule keys process selection on the first term alone \cite{bezos2015}. ``Architectural'' is not
a category of topics. It is the high-product region of this calculus, which is why teams
adopting decision records report that their hardest problem is deciding what counts as
architectural \cite{weiss2024}, and why adoption in the wild is shallow, with half of adopting
repositories abandoning after fewer than five records \cite{buchgeher2023}. \emph{Evidence:}
decision graph plus git history; partially inferable, confirmable by a human.
\end{axis}

\begin{axis}[Commitment state and restriction force]\label{ax:state}
Kruchten's ordered promotion scale, from idea to tentative to decided to approved, with
challenged, rejected, and obsolesced as demotions, with consistency propagation along the
decision graph's typed edges: constrains, forbids, enables, subsumes, comprises, conflicts-with,
overrides, is-alternative-to \cite{kruchten2004}. And, per decision, its restriction force: how
strongly its rules and constraints prune the future option space. A ban prunes maximally and
leaves zero code residue \cite{kruchten2004}. A convention prunes weakly. Jansen and Bosch's
deltas make the composition explicit: an architecture at time $t$ is a fold over the decision
sequence \cite{jansenbosch2005}. QOC said it first and plainest: a design is a point in a space
constructed by questions, which are dimensions, options, which are values, and criteria, which
are the value field \cite{maclean1991}. \emph{Evidence:} requires the decision store,
historically vaporized; see \S\ref{sec:agent} for why an agent dissolves the four documented
causes of that vaporization.
\end{axis}

One honest note closes the tier. The calculus describes the decision, not the decider. The
empirical decision literature adds decider variables: the thin spread of domain knowledge across
a team \cite{curtis1988}; problem structuredness selecting between rational comparison and
recognition-primed deciding \cite{zannier2007}; bounded rationality and its documented biases
\cite{vanvliettang2016}. These govern \emph{how} moves get made, and they mark where an agent's
estimates must substitute for records that were never going to be written.

\subsection{Tier IV: the value field and the dynamics}\label{sec:dynamics}

The last tier separates two things the field habitually merges with the coordinates. The first
is how good a position is for a given purpose, which is a field over the space and changes when
goals change. The second is which way the terrain slopes when nobody is pushing, which is a
dynamics and does not care about goals at all.

\begin{axis}[The value field]\label{ax:value}
Quality attributes form a field \emph{over} the space, not axes \emph{of} it: Fielding's
thirteen properties, ISO 25010's eight characteristics, and the quality-attribute-scenario
formulation that makes each one measurable \cite{fielding2000,iso25010,bass2012}. Fielding's
Table 3--6 is a sampled gradient, each constraint's signed property increments with
non-additivity explicitly flagged. ATAM's sensitivity points are parameters with large
$\partial q / \partial x$, and its tradeoff points are parameters where two attributes'
derivatives have opposite sign \cite{atam2000}. A project's criteria weights, whether the
utility tree, QOC's criteria, or MADR's decision drivers \cite{madr2018}, are its direction of
ascent. Keeping the field separate from the coordinates is what lets goals change without
remapping the terrain. And it is Lane's structure exactly: functional dimensions from
requirements, correlated by learned rules to structural dimensions of choice, the rules being
the empirical map between the two spaces \cite{lane1990}.
\end{axis}

\begin{axis}[Entropic drift and the maintenance ratio]\label{ax:drift}
Lehman's laws supply the vector field: under E-type change, structure degrades unless
anti-regressive work is spent, and the controllable variable is the allocation between
progressive and anti-regressive effort \cite{lehman1980}. The empirical decay literature
confirms drift as the default trajectory across hundreds of open-source versions
\cite{behnamghader2017}, and the smell and hotspot catalogs \cite{azadi2019,mo2016} name the
attractor's neighborhoods. This is where the one-way axes live. Cyclic cores, hub overloads,
and modularity violations, meaning co-change without structural dependency, the misalignment of
latent and manifest coupling that is measurable only from history \cite{mo2016}, are potentials
rather than coordinates. They are descended into by accident and climbed out of only by
spending energy.
\end{axis}

\begin{axis}[Framing: the attention vector]\label{ax:framing}
The design-cognition literature offers a hard datum. Problem and solution spaces co-evolve
\cite{dorstcross2001}. The frame determines which dimensions are salient \cite{schon1983}.
Design is constraint-dominated, with hard and soft constraints, leaky decompositions, and
personalized stopping rules \cite{goelpirolli1992}. The space is not experienced as fixed. The
defensible reconciliation, argued at length in \S\ref{sec:limits}, is a fixed superset of
generators plus a time-varying salience weighting. Framing rotates the active subspace; it does
not create axes. G\"ardenfors' machinery contains exactly this knob, context-dependent dimension
weights \cite{gardenfors2000}, and Gero's FBS ontology gives the tier structure that localizes
reframing as explicit operators, reformulation types 1 through 3, between function, behaviour,
and structure spaces \cite{gero2004}.
\end{axis}

\subsection{The basis at a glance}

\begin{center}
\small
\begin{longtable}{@{}p{0.26\textwidth}p{0.30\textwidth}p{0.20\textwidth}p{0.16\textwidth}@{}}
\toprule
\textbf{Variable} & \textbf{Formal structure} & \textbf{Key convergences} & \textbf{Evidence} \\
\midrule
\endhead
Cut placement $\Pi_i$, $\varepsilon_i$ (4 graphs) & partition lattice; near-block-diagonal matrices & Simon; Alexander; SMC; clustering objectives & repo / traces / org / domain \\
Granularity (per graph) & scale ladder; basic level & Rosch; Chromium and Istio two-way traffic & repo, manifests \\
Congruences $\kappa_{OC},\kappa_{PC},\kappa_{CR}$ & matrix congruence & Conway; mirroring studies; DDD; migration post-mortems & repo + git + org \\
Boundary interaction type & product: topology $\times$ sync $\times$ mode $\times$ guarantees $\times$ state locus & boxology; connector taxonomy; coupling ladder & repo \\
Interface thinness, uniformity & bits leaked; cut bandwidth & Parnas; Simon; Baldwin; Allen and Khoshgoftaar & repo \\
Uses order & partial order; ideals = subsets & Parnas 1979; Simon's stable forms & repo \\
Abstraction order & chain of state-space quotients & Dijkstra; refinement calculus & repo + human \\
Dominance order / design rules & ACN; DRH; refinement orders & Baldwin and Clark; Cai and Sullivan; Martin; contracts & repo \\
Dependency orientation & vector-field constraint on $(I,A)$ & Martin; hexagonal/onion/clean & repo \\
Reachability density & transitive-closure density; DL & MacCormack $\equiv$ Lakos; Mo et al. & repo \\
Volatility field & weighted decision set; $\sigma_i$ & Parnas 1972; options theory & git + human \\
Binding time & ordered scale per decision & 5 independent frameworks & repo \\
Reversal cost / blast radius / dependents & product over decision graph & Fowler; Kruchten; options; Bezos & decisions + git \\
Commitment state, restriction force & ordered states; typed decision edges; deltas & Kruchten; Jansen and Bosch; QOC & decision store \\
Value field, criteria weights & gradient field; utility tree & Fielding; ATAM; Lane; ISO 25010 & human + measurement \\
Entropic drift, maintenance ratio & vector field; potentials & Lehman; decay studies; smells & git \\
Framing / attention vector & salience weighting over fixed generators & Schön; Dorst and Cross; Gärdenfors; FBS & interaction \\
\bottomrule
\end{longtable}
\end{center}

Seventeen entries. But the count overstates the independence, because the deep structure is the
four-tier claim: cuts and their congruences; orders over the cuts; a commitment calculus over
decisions; a field and a flow over the whole. Everything in the literature we surveyed, every
metric, style, tactic, smell, and principle, reduces to a position, region, move, gradient, or
potential expressed in these terms. \S\ref{sec:reduction} performs the reductions explicitly.

% ===========================================================================
\section{The space: geometry, prototypes, moves, metric}\label{sec:space}
% ===========================================================================

A list of variables is not yet a space. For what more is needed, we borrow from cognitive
science. G\"ardenfors' conceptual spaces are a simple, well-studied picture of how minds hold
concepts: qualities become dimensions, concepts become regions, and similarity becomes nearness
\cite{gardenfors2000}. A reader meeting the idea for the first time needs only that much. The
theory then specifies what must be checked rather than assumed: that the dimensions have
intrinsic structure; that similarity decreases lawfully with distance \cite{shepard1987}; that
betweenness makes sense; that natural regions are convex and organized around prototypes,
which is Voronoi structure; and that dimensions divide in a principled way into integral and
separable bundles. We take these requirements in turn, honestly.

\subsection{Coordinates and locality}

The dimensions of \S\ref{sec:basis} are of mixed type by construction. There are ordinal scales:
binding time, coupling ladders, commitment states. There are continuous scalars: $\varepsilon$,
reachability density, the congruences. There are categorical dimensions with structure: the
interaction type as a product of small enumerations. And there are full partial orders: uses,
abstraction, dominance. The space is therefore a hybrid, a discrete relational skeleton of
graphs and orders woven with continuous coordinates of density, congruence, and semantic
similarity. Betweenness holds on the continuous sub-spaces. Between synchronous and asynchronous
sits deferred-synchronous, and the connector taxonomy enumerates exactly such intermediates
\cite{mehta2000}. Betweenness fails, as expected, on the discrete skeleton, where paths replace
interpolation. This is not a defect. It is the same hybrid structure that navigation research
already handles by treating betweenness as optional on graph-like domains \cite{behrens2018}.

Integral versus separable bundles can be read off the empirical record. The performance bundle
of latency, throughput, and resource cost, and the trade surface of consistency and
availability, behave as integral: one cannot be valued without valuing the others, and they
combine in a Euclidean-like way in engineers' trade-off talk. The tier-one structural axes are
largely separable. The paper's central empirical finding (\S\ref{sec:empirical}) is precisely
that the logical cut and the physical cut, which the field habitually fuses, vary independently.

\subsection{Styles and patterns are the prototype system}\label{sec:prototypes}

The G\"ardenfors reading of the styles literature is exact, not analogical. A style, in
Fielding's definition, is a coordinated set of constraints \cite{fielding2000}. That makes it a
\emph{region}: the set of configurations satisfying the constraints, with the null style, the
empty constraint set, as the whole space, and constraint addition as region refinement.
Boxology's tables assign each named style its coordinates on the explicit dimension set
\cite{boxology1997}. Our third research lane performs the assignment for pipes-and-filters in
full: passed-stream data mode entails no shared state, which permits asynchronous incremental
processing, which yields the functional-composition proof rule and the reuse properties. The
region is coherent, a natural concept rather than an arbitrary conjunction, because its
coordinate values are \emph{mutually entailing}. That is what makes a style learnable from few
examples and nameable in one word. It is exactly the learnability-and-communication pressure
that G\"ardenfors' convexity criterion encodes \cite{gardenfors2000,jager2007}.

Patterns are prototypes at finer grain. Alexander's own arc ran from coordinates, the 1964
misfit-graph program he later repudiated, to prototypes with forces in 1977, and the software
patterns movement imported the prototype form with the forces sections intact
\cite{alexander1977,gamma1994}. Patterns are the community's informally collected basis vectors,
each one a named equilibrium of forces with a region of variation around it, ``without ever
doing it the same way twice.'' The recognition literature confirms that this is how experts
actually operate. Recognition-primed decisions land in a prototype cell and read off the
associated action \cite{klein1998}. Programming plans and beacons are the code-level instance
\cite{soloway1984}. And the recovered structure of real systems clusters. Most large codebases
sit in the core-periphery region \cite{baldwin2014}. Martin's main sequence is a prototype line
with named pathological corners \cite{martin2002}. Prototype structure is not imposed by the
theory. It keeps being found.

Two honest qualifications. Styles show family resemblance: no necessary-and-sufficient
conditions, graded membership, ``sort of microservices.'' So convexity must be read as a
pressure rather than a law, and genuinely disjunctive styles violate it \cite{hernandez2017}.
And expert categorization is part exemplar, the war stories, ``like the system at Corp X,'' and
part causal theory, classification by \emph{why}. Rationale is the scarcest, most-sought
information in comprehension studies \cite{latoza2010}. A pure geometric prototype store is
necessary but not sufficient. The space needs its instances and its reasons attached.

\subsection{The move set}\label{sec:moves}

The space's operations are already cataloged, at three grains, by literatures that do not cite
each other:

\begin{itemize}[leftmargin=1.6em]
\item \textbf{Generators.} Baldwin and Clark's six modular operators: split, substitute,
augment, exclude, invert, port \cite{baldwinclark2000}. Every entry in the refactoring catalogs
is an instance. \emph{Invert}, which extracts a repeated function and raises it into shared
infrastructure, is the origin of every platform layer.
\item \textbf{Unit moves along single axes.} The tactics catalogs, several of which name their
axis outright: defer binding (Axis~\ref{ax:binding}); use an intermediary, which is
intermediation depth; introduce concurrency, which is synchronicity; maintain multiple copies,
which is replication \cite{bass2012}.
\item \textbf{Region jumps.} Style adoption as a correlated multi-dimensional displacement, with
Fielding's REST derivation as the audited exemplar: six constraint-moves from the null style,
each logged with its induced property increments and costs \cite{fielding2000}.
\end{itemize}

A decision, in the commitment calculus, is a move plus a restriction. It displaces the point and
prunes the reachable region through design rules and constraints \cite{jansenbosch2005}. The
typed edges of the decision graph \cite{kruchten2004} are the algebra by which moves compose,
conflict, enable, and override.

\subsection{The metric: asymmetric, use-shaped, and answerable to Tversky}\label{sec:metric}

The most consequential design choice in the whole construction is the metric, and this is a
point where it becomes important to slow down, because the cognitive literature issues a warning
here that the architecture literature has never processed. Tversky's classical result is that
human similarity violates every metric axiom. It is asymmetric, since the variant is compared to
the prototype and not conversely. It fails the triangle inequality. It is context-dependent
\cite{tversky1977}. Any theory that installs a fixed symmetric Euclidean metric over
architectural descriptions is empirically false about the judgments it must model.

The escape is already present in the material. Distance in a decision space is re-decision
cost: how expensive it is to move this system from here to there, under the commitment
calculus. And that quantity is inherently asymmetric. Splitting a monolith and merging
microservices are not inverse operations at equal price. Adding an edge is cheap and removing
it is dear; the one-way axes of \S\ref{sec:dynamics} are the extreme cases. It is also
inherently context-dependent, since the cost depends on binding times, accrued dependents, and
the current frame. The neurally grounded form of exactly this metric exists. It is the
successor representation, under which nearness is expected future co-occupancy under the
agent's policy: predictive, directional, graph-friendly, and empirically supported as the
hippocampal map's actual code \cite{dayan1993,stachenfeld2017}. Seed the metric with surface
structure, then let realized navigation reshape it, watching which decisions actually follow
which and which units actually co-change. The co-change matrix is precisely an empirical
co-navigation record. Tversky's asymmetry then stops being an objection and becomes a
prediction. Architectural similarity judgments \emph{should} be asymmetric, in the direction of
re-decision cost, because the space's true metric is reachability under work.

This yields the paper's second structural claim:

\begin{claim}\label{claim:metric}
The conceptual space of architecture is a hybrid manifold with an asymmetric, use-shaped metric.
Its distances are re-decision costs under the commitment calculus, seeded by structural
similarity and reshaped by realized co-change. Its value field, held separate, orients
navigation without deforming the terrain.
\end{claim}

% ===========================================================================
\section{Reduction tests: the familiar vocabulary is composite}\label{sec:reduction}
% ===========================================================================

A basis earns its title by rewriting the field's vocabulary as compositions. We perform the six
most load-bearing reductions. Each one dissolves a standing confusion.

\subsection{``Microservices'' is four coordinates, not one}

The monolith-versus-microservices axis, as debated, conflates four different choices
\cite{fowler2014,newman2015}. One is deployment granularity: how many independently deployable
units there are. Another is failure isolation: whether a call stays in one process or crosses a
network. A third is data ownership: one transactional store, or a store per service with
eventual consistency between them. The fourth is organizational alignment: whether services are
sized to teams, which is the $\kappa_{OC}$ term. The empirical record separates them cleanly. Prime Video kept its
component decomposition and deleted only the network boundary \cite{primevideo2023}. Shopify
added enforced module boundaries while keeping one deployable \cite{shopify2019}. Uber added
coarse logical domains over fine physical services \cite{uberdoma2020}. ``Modular monolith'' is
not a compromise. It is a coordinate assignment: logical cut fine, physical cut coarse. The
debate dissolves into force-weighting on four separable axes.

\subsection{``Layering'' is three partial orders}

Containment, which is Simon's boxes-in-boxes and the module guide's is-submodule-of
\cite{parnas1985}; use, from Parnas 1979; and abstraction, from Dijkstra 1968, are independent
orders that the word ``layer'' merges. The classical corpus itself insists on the distinction
\cite{parnas1972}. Systems exist with clean containment and cyclic use, and with layered use
that abstracts nothing. Conformance rules in practice, such as ``controllers may not access
repositories directly,'' are constraints on the use order. C4's levels are resolutions of the
containment order \cite{c4model}. Ports-and-adapters is a constraint tying the use order to the
abstraction gradient (Axis~\ref{ax:orientation}).

\subsection{``Coupling'' is a product, plus a residual, plus a ghost}

Structural coupling is the boundary interaction type (Axis~\ref{sec:boundarytype}), itself a
product rather than a scalar. Logical coupling is a different measurement entirely: co-change in
history. And the \emph{misalignment} of the two, co-change without structural dependency,
correlates with error-proneness and change-proneness across many projects. It is a violated
information-hiding boundary leaving a historical ghost \cite{mo2016}. Cohesion, the traditional
partner concept, may not be a dimension at all. Its metrics fail measurement-theoretic axioms
and disagree with each other \cite{briand1998}, while the information-theoretic view renders it
as the residual of the cut, what remains uncompressed when the partition is chosen
\cite{allen1999,andritsos2005}. We therefore admit coupling-type, logical coupling, and their
misalignment into the basis, and demote cohesion to a derived quantity. That is a specific,
falsifiable ontological claim.

\subsection{``Hard to change'' is a computable product}

Fowler's definition-with-a-wink \cite{fowler2003} decomposes as Axis~\ref{ax:reversal}:
reversal cost times blast radius times dependents accrued, modulated by binding time. Each term
is separately estimable, and each is separately manipulable. The point of the decomposition is
practical. An agent cannot act on ``perceived as hard to change.'' It can compute accruing
dependents, measure blast radius through the visibility matrix, and flag binding times that
arrived earlier than necessary.

\subsection{``The architecture'' recovered from code is a projection, and the recovery
literature proves it}

Recovery algorithms each optimize a formalized theory of what a module is: dependency density in
Bunch's MQ \cite{mancoridis1999}, recurring structural idiom in ACDC \cite{tzerpos2000}, concern
coherence in ARC's topic models \cite{garcia2011}, information equivalence in LIMBO
\cite{andritsos2005}. Measured against engineer-certified ground truths, all score poorly, and
they disagree with each other roughly as much as with the truth \cite{garcia2013,lutellier2015}.
On the single-partition theory of architecture this is an embarrassment. On the present theory
it is a \emph{measurement}. The objectives capture different real dimensions. The engineer's
decomposition integrates variables the dependency graph does not carry: ownership, deployment,
domain semantics, history. Any one partition is a lossy projection of a multi-dimensional latent
state, and the disagreement pattern across objectives is itself a fingerprint of where a system
sits in the space.

\subsection{``Separation of concerns'' and the tyranny of one cut}

A concern is a projection of the design space onto a subspace, which is Dijkstra's own
formulation \cite{dijkstra1974}. The multi-view doctrine, 4+1 and the viewpoint standards,
institutionalizes a small atlas of such projections keyed to stakeholder concerns
\cite{kruchten1995}. The file tree realizes exactly one cut, and whatever does not align with
the chosen dominant dimension becomes scattered and tangled. This is the tyranny of the dominant
decomposition \cite{tarr1999}. Scattering is measurable, and scattered concerns carry
significantly more defects \cite{eaddy2008}. The latent state therefore includes \emph{which}
dimension was chosen dominant, the concern-to-code distribution, a multi-label field over units,
and the misalignment between them. None of these can be expressed by any file tree. All of them
are estimable (\S\ref{sec:agent}).

% ===========================================================================
\section{The empirical shape of the space}\label{sec:empirical}
% ===========================================================================

Theory proposes; the deployed world disposes. Our sixth research lane surveyed eighteen major
systems through their own architects' accounts, the AOSA corpus and primary engineering sources,
along with seven canonical public debates, five industrial macro-migrations, and the
mining-software-repositories literature. Three findings shape the theory.

\subsection{The decision points localize the dimensions}

When both sides of a real dispute articulate their forces, the disagreement names an axis.
Tanenbaum and Torvalds argued isolation-boundary placement and its costs, IPC overhead against
address-space protection \cite{tanenbaum1992}. The same axis resurfaced, thirty years later and
in the participants' own terms, in Postgres's process-versus-threads deliberations
\cite{postgres2023} and in every browser's multi-process migration \cite{chromium2009}. Hipp's
``SQLite competes with fopen()'' collapses a dozen surface differences into one coordinate:
where the process, trust, and administration boundary sits between application and engine
\cite{sqlitearch}. React Fiber's ``a fiber is a reimplementation of the stack'' names scheduling
ownership: implicit in the language runtime, or reified as data \cite{fiber2016}. The systemd
wars conflated two axes, declarative versus imperative configuration and integrated platform
versus orthogonal components, and the temperature of that debate fell where the axes were
distinguished \cite{poettering2010}. In every case the engineers' own vocabulary, not the
analyst's, supplies the dimension. The folk coordinates and the theoretical ones agree.

\subsection{Two-way traffic identifies the free coordinates}

Documented bidirectional migration under articulated forces is the strongest available evidence
that an axis is a genuine degree of freedom. It exists for isolation granularity: Chromium
finer, Istio coarser. For logical-physical congruence: Prime Video decoupled them one way, Uber
the other. For concurrency model: Apache and nginx converging from opposite poles. For
centralization: ZeroMQ removed the broker RabbitMQ kept; Git decentralized what SVN
centralized; ZooKeeper deliberately re-centralized coordination \cite{hunt2010}. For binding
time: nginx compile-time against Eclipse hot-install. And for state representation: snapshots
versus deltas versus log-as-truth, with Git famously choosing snapshots at the logical level
and deltas at the physical level, splitting the axis across an abstraction boundary
\cite{aosagit,aosazeromq,segment2018}. By contrast, the one-way axes, cyclic-core growth, hub
overload, and scattering, attract only unintentional traffic inward and deliberate traffic
outward. They are potentials, per \S\ref{sec:dynamics}.

\subsection{What varies together, and what refuses to}

Read the eighteen-system variation table columnwise. The columns are process model, concurrency
model, extension mechanism, topology, state management, isolation boundary, binding, and
synchrony. Every one of them shows multiple values across respected systems, and the columns
vary with substantial independence. That is the empirical signature of separable dimensions. The strongest single
regularity is negative. Attempts to force one column to determine another, deployment
granularity dictating module granularity, or org structure ignored while cutting code, are the
named root causes in the migration post-mortems \cite{primevideo2023,shopify2019,uberdoma2020}.
Where the population does concentrate, the concentrations are the prototype regions and role
structures the geometry predicts (\S\ref{sec:prototypes}). Most large systems are
core-periphery \cite{baldwin2014}. Production call graphs are heavy-tailed, with a few hub
services carrying enormous fan-in \cite{luo2021}. Fan-in is power-law distributed across every
code ecosystem measured \cite{louridas2008}. This includes the hub role that pure
cohesion-and-coupling theories systematically misjudge. Utilities are coupled to everyone.
Quarantining them is an idiom every practitioner knows, and only the pattern-based recovery
objective encodes it \cite{tzerpos2000}.

% ===========================================================================
\section{What a machine can measure}\label{sec:agent}
% ===========================================================================

The computability test now pays out. We assess, variable by variable, what suffices to estimate
the basis, and we measure the gap between that and what coding agents actually represent today.

\subsection{The evidence ladder}

Start with a \textbf{bare repository} and unglamorous tooling: a dependency extractor,
Tarjan's algorithm, matrix powers, tree-sitter, an embedding model. From that alone an agent
can already estimate a great deal. It can compute the code cut and its candidate alternatives
under all four recovery objectives, whose disagreement profile is itself a fingerprint. It can
compute reachability density, decoupling level, core size, element roles, cycles, levelization,
and Martin's $(I,A,D)$. It can read boundary interaction types out of interface signatures and
connector idioms, along with binding times, declared visibility, and API surfaces. And it can
infer technology bindings per module, naming-convention stereotypes, concern vectors from
identifiers and comments, and the dominant-decomposition choice, meaning what the tree is
organized by. One methodological caveat travels with all of it: extraction quality changes
results as much as algorithm choice \cite{lutellier2015}.

From \textbf{git history}: churn and hotspots; logical coupling; modularity violations, the
highest-value single derived signal \cite{mo2016}; interface instability; ownership and
$\kappa_{OC}$; and drift dating, meaning when a divergence appeared and in which commit, which
often links to \emph{why}. From \textbf{runtime traces}: the realized interaction graph as
against the potential one; hot paths; dynamic feature location. Requiring \textbf{human
intent}, irreducibly but cheaply: the intended module map and allowed-dependency relation,
which are the reflexion inputs \cite{murphy1995}, the fitness priorities, and rationale for
tolerated divergences.

Of the seventeen basis entries, roughly nine are estimable from the bare repo, four more fall
out of history, and only the intent-carriers need a person, via the draft-correct-enforce loop
that reflexion proved and that language models make cheap.

\subsection{The gap, and the tell}

Against this, the composite working representation of current agents, across the major coding
assistants and open scaffolds of 2024 through 2026, is: the file tree, a symbol and reference
graph, lexical and embedding retrieval, and an optional human-written prose file of guidance
\cite{aider2023,repograph2024}. Hold that up against the basis. No intended-module map. No
allowed-dependency relation. No concern vectors. No propagation structure. No history
aggregation. No roles. No commitment calculus. No value weights. The agents hold the
observables and none of the latents.

The prose context file is the tell, and it deserves to be named as such. Teams hand-maintain
\texttt{CLAUDE.md} and its cousins because they know the latents matter. The file is a
human-powered cache of machine-computable quantities, lossy and stale by construction. It is
the architecture-description pathology, dual maintenance, no incremental payoff, semantic
distance from code, that killed the ADLs \cite{malavolta2013}, replayed in miniature. The
survivors of that extinction event are instructive. Architecture description persisted only
where it became \emph{derived from} or \emph{enforced upon} the code: build-system visibility
rules, typed interface schemas, conformance tests, deployment manifests \cite{fitness2017}. The
design implication for an agent is direct. Represent the basis as derived, continuously
refreshed quantities plus a small store of confirmed intent, and let no fact live only in
prose.

\subsection{Why the decision store becomes viable now}

The commitment calculus (\S\ref{sec:commitment}) requires the one thing history never had: a
populated decision store. Rationale capture failed for four documented, mutually reinforcing
human-economic reasons. The capturer is not the beneficiary \cite{grudin1996}. Formalization
carries cognitive overhead and demands premature structure \cite{shipman1999}. Capture intrudes
on design flow mid-flight. And practitioners perceive the value as deferred and diffuse
\cite{tang2006}. Notice what these four have in common. Every one is an asymmetry of human
cost, not a defect of the representations. And every one dissolves for an agent that makes
decisions itself, records them as a zero-marginal-cost side effect at the moment of deciding,
infers structure rather than demanding it, and is its own future consumer. The mined reality
that most real decisions live in issues, pull requests, and chat rather than in decision
records \cite{bhat2017} becomes, for such an agent, a recoverable backlog rather than a loss.
And the one strong empirical result from the LLM-generation literature, that a recency window
of three to five prior decisions from the same project is the best conditioning signal for the
next one \cite{contextmatters2026}, is exactly what the trajectory theory predicts. The signal
is in the path, and current representations carry it only by accident.

% ===========================================================================
\section{Honest limits and open problems}\label{sec:limits}
% ===========================================================================

Four tensions survive this synthesis undischarged. We state them as constraints on the theory's
ambitions, and two of them convert into its sharpest research questions.

\textbf{The fixed basis versus co-evolution.} The design-cognition record is unambiguous.
Problem and solution spaces co-evolve \cite{dorstcross2001}. Framing chooses what counts as a
dimension \cite{schon1983}. Shaw and Brooks insist the space evolves mid-design
\cite{shaw2012,brooks2010}. Our reconciliation, a fixed superset of generators with a
time-varying attention weighting, is defensible. The basis variables recur across sixty years
and eighteen systems, and that stability is evidence. But it is a hypothesis, not a theorem.
Its falsifiable content is this: protocol studies of architects should find reframings that
re-weight the basis, and should not find recurring decision dimensions that resist expression
in it. A reframing that genuinely manufactures a new generator would falsify the fixed-superset
claim.

\textbf{The identification problem.} Interaction minimization and change localization are
different objectives that coincide only if change propagates uniformly along interaction edges
(\S\ref{sec:object}). No one has measured the actual kernel, the conditional probability that a
change at $u$ implicates $v$, as a function of the edge types and orders between them. The
co-change record makes it measurable for the first time at scale. This is the paper's most
concrete open empirical problem, and its answer calibrates every cut-quality metric in use.

\textbf{The metric's empirical status.} Claim~\ref{claim:metric} is theory-consistent. The
asymmetry matches Tversky, and the successor-representation structure matches the navigation
literature. But it is unvalidated for architecture. No one has collected architects' similarity
and substitutability judgments at scale and tested whether they embed with the predicted
asymmetries, and whether generalization decays with the embedded distance \cite{shepard1987}.
Until that experiment is run, the metric is the best-motivated candidate, not a measured fact.

\textbf{Convexity at the edges.} Family-resemblance styles, disjunctive categories, and the
hybrid discrete-continuous structure mean the Voronoi picture applies cleanly only region by
region. We have treated convexity as a pressure shaping nameable vocabulary, which is the
defensible reading \cite{jager2007,hernandez2017}. But it demotes one of the theory's prettiest
formal properties to a tendency, and we would rather say so plainly than pretend otherwise.

% ===========================================================================
\section{Falsifiable predictions}
% ===========================================================================

\begin{pred}
\emph{Asymmetric similarity.} Architects' pairwise judgments of system similarity and migration
difficulty will be systematically asymmetric, with the asymmetry oriented by re-decision cost,
meaning dependents accrued and binding time, not by surface feature counts.
\end{pred}

\begin{pred}
\emph{The congruence terms dominate.} In matched samples, $\kappa_{OC}$ and the
logical-physical congruence $\kappa_{CR}$ will predict maintenance outcomes better than any
single-graph structural metric, including propagation cost, extending the Vista result
\cite{nagappan2008} from defects to velocity.
\end{pred}

\begin{pred}
\emph{The two-way/one-way partition is stable.} New industrial migrations will continue to move
bidirectionally on the coordinate axes, granularity, centralization, binding, concurrency
model, state representation, and unidirectionally on the potentials, cyclic core, scattering,
hub overload. An intentional, force-justified migration \emph{toward} a larger cyclic core
would count against the theory.
\end{pred}

\begin{pred}
\emph{Recovery disagreement is lawful.} The disagreement profile among recovery objectives,
density, idiom, concern, information, on a given system is reproducible and predicts which
conformance rules its maintainers will accept, because both are functions of the same latent
position.
\end{pred}

\begin{pred}
\emph{An agent with the basis beats an agent with the observables.} On repository-level change
tasks, an agent whose context is the estimated basis, meaning the intended map, the
allowed-dependency relation, propagation structure, hotspots, roles, concern vectors, and the
recent decision trajectory, will outperform an agent with the file tree, symbol graph, and
embeddings at equal token budget. The margin will be largest on tasks whose blast radius
crosses module boundaries.
\end{pred}

\begin{pred}
\emph{Decision-trajectory signal.} Conditioning next-decision prediction on the typed decision
graph, its states, edges, and restriction force, will beat the flat recency window of
\cite{contextmatters2026}, because the window is a degenerate projection of the graph.
\end{pred}

% ===========================================================================
\section{Coda: building the space}
% ===========================================================================

The construction recipe follows the theory. Coordinates: place every code unit by its
structural address, graph position, role, concern vector, and binding profile, and place every
decision by its commitment calculus. Locality: neighborhoods from dependency, call, co-change,
and semantic distance, with betweenness only where the sub-space is continuous. Prototypes: the
styles, patterns, and role archetypes as landmarks, and a repository's own core abstractions
and canonical instances as its local landmarks. Operations: the six operators, the tactics, and
the region jumps, exposed as the only mutation vocabulary. Metric: seeded by structure,
reshaped by realized co-navigation, asymmetric by construction. Value field: the project's
criteria weights, held separate, orienting navigation without deforming terrain.

An agent so equipped does not read a codebase the way a search engine reads a corpus. It knows
where it is, what is near, which moves are legal, what each move costs, and which way is up.
That, and not a longer context window, and not a better prose file, is what it means for a
machine to understand architecture.

The deeper wager of this paper is methodological. The field's sixty-year literature, read as a
single corpus, is not a pile of opinions. It is a set of independent measurements of one
underlying object, taken with different instruments. When the measurements converge, and we
have shown that they converge, from Parnas's decision trees to Baldwin's options to boxology's
tables to the migration post-mortems, the object is real. The latent variables of software
architecture are not a metaphor awaiting rigor. They are the oldest finding in the field,
waiting to be assembled.

\begin{thebibliography}{99}
\small

\bibitem{parnas1972} Parnas, D.L. (1972). On the criteria to be used in decomposing systems into
modules. \emph{Communications of the ACM} 15(12):1053--1058.

\bibitem{parnas1976} Parnas, D.L. (1976). On the design and development of program families.
\emph{IEEE Transactions on Software Engineering} SE-2(1):1--9.

\bibitem{parnas1979} Parnas, D.L. (1979). Designing software for ease of extension and
contraction. \emph{IEEE TSE} SE-5(2):128--138.

\bibitem{parnas1985} Parnas, D.L., Clements, P.C., Weiss, D.M. (1985). The modular structure of
complex systems. \emph{IEEE TSE} SE-11(3):259--266.

\bibitem{dijkstra1968} Dijkstra, E.W. (1968). The structure of the ``THE''-multiprogramming
system. \emph{Communications of the ACM} 11(5):341--346.

\bibitem{dijkstra1974} Dijkstra, E.W. (1974). On the role of scientific thought. EWD~447.

\bibitem{simon1962} Simon, H.A. (1962). The architecture of complexity. \emph{Proceedings of the
American Philosophical Society} 106(6):467--482.

\bibitem{simonando1961} Simon, H.A., Ando, A. (1961). Aggregation of variables in dynamic
systems. \emph{Econometrica} 29:111--138.

\bibitem{conway1968} Conway, M.E. (1968). How do committees invent? \emph{Datamation}
14(4):28--31.

\bibitem{brooks2010} Brooks, F.P. (2010). \emph{The Design of Design}. Addison-Wesley.

\bibitem{stevens1974} Stevens, W.P., Myers, G.J., Constantine, L.L. (1974). Structured design.
\emph{IBM Systems Journal} 13(2):115--139.

\bibitem{lehman1980} Lehman, M.M. (1980). Programs, life cycles, and laws of software evolution.
\emph{Proceedings of the IEEE} 68(9):1060--1076.

\bibitem{alexander1964} Alexander, C. (1964). \emph{Notes on the Synthesis of Form}. Harvard
University Press.

\bibitem{alexander1977} Alexander, C., et al. (1977). \emph{A Pattern Language}. Oxford
University Press.

\bibitem{baldwinclark2000} Baldwin, C.Y., Clark, K.B. (2000). \emph{Design Rules, Vol.~1: The
Power of Modularity}. MIT Press.

\bibitem{baldwin2008} Baldwin, C.Y. (2008). Where do transactions come from? Modularity,
transactions, and the boundaries of firms. \emph{Industrial and Corporate Change} 17(1):155--195.

\bibitem{colferbaldwin2016} Colfer, L.J., Baldwin, C.Y. (2016). The mirroring hypothesis: theory,
evidence, and exceptions. \emph{Industrial and Corporate Change} 25(5):709--738.

\bibitem{steward1981} Steward, D.V. (1981). The design structure system. \emph{IEEE Transactions
on Engineering Management} EM-28(3):71--74.

\bibitem{maccormack2006} MacCormack, A., Rusnak, J., Baldwin, C.Y. (2006). Exploring the
structure of complex software designs: an empirical study of open source and proprietary code.
\emph{Management Science} 52(7):1015--1030.

\bibitem{maccormack2012} MacCormack, A., Baldwin, C.Y., Rusnak, J. (2012). Exploring the duality
between product and organizational architectures: a test of the ``mirroring'' hypothesis.
\emph{Research Policy} 41(8):1309--1324.

\bibitem{baldwin2014} Baldwin, C.Y., MacCormack, A., Rusnak, J. (2014). Hidden structure: using
network methods to map system architecture. \emph{Research Policy} 43(8):1381--1397.

\bibitem{nagappan2008} Nagappan, N., Murphy, B., Basili, V. (2008). The influence of
organizational structure on software quality. \emph{Proc.\ ICSE 2008}.

\bibitem{cataldo2008} Cataldo, M., Herbsleb, J.D., Carley, K.M. (2008). Socio-technical
congruence. \emph{Proc.\ ESEM 2008}.

\bibitem{bird2011} Bird, C., Nagappan, N., Murphy, B., Gall, H., Devanbu, P. (2011). Don't touch
my code! Examining the effects of ownership on software quality. \emph{Proc.\ ESEC/FSE 2011}.

\bibitem{accelerate2018} Forsgren, N., Humble, J., Kim, G. (2018). \emph{Accelerate}. IT
Revolution Press.

\bibitem{sullivan2001} Sullivan, K.J., Griswold, W.G., Cai, Y., Hallen, B. (2001). The structure
and value of modularity in software design. \emph{Proc.\ ESEC/FSE 2001}.

\bibitem{cai2006} Cai, Y., Sullivan, K.J. (2006). Modularity analysis of logical design models.
\emph{Proc.\ ASE 2006}.

\bibitem{mo2016} Mo, R., Cai, Y., Kazman, R., Xiao, L., Feng, Q. (2016). Decoupling level: a new
metric for architectural maintenance complexity. \emph{Proc.\ ICSE 2016}.

\bibitem{lakos1996} Lakos, J. (1996). \emph{Large-Scale C++ Software Design}. Addison-Wesley.

\bibitem{martin2002} Martin, R.C. (2002). \emph{Agile Software Development: Principles,
Patterns, and Practices}. Prentice Hall.

\bibitem{allen1999} Allen, E.B., Khoshgoftaar, T.M. (1999). Measuring coupling and cohesion: an
information-theory approach. \emph{Proc.\ IEEE METRICS 1999}.

\bibitem{briand1998} Briand, L.C., Daly, J.W., Wüst, J. (1998--99). A unified framework for
coupling/cohesion measurement in object-oriented systems. \emph{IEEE TSE} 25(1) / \emph{ESE} 3(1).

\bibitem{mancoridis1999} Mancoridis, S., Mitchell, B.S., et al. (1999). Bunch: a clustering tool
for the recovery and maintenance of software system structures. \emph{Proc.\ ICSM 1999}.

\bibitem{andritsos2005} Andritsos, P., Tzerpos, V. (2005). Information-theoretic software
clustering. \emph{IEEE TSE} 31(2):150--165.

\bibitem{dealfaro2001} de Alfaro, L., Henzinger, T.A. (2001). Interface automata. \emph{Proc.\
ESEC/FSE 2001}.

\bibitem{benveniste2018} Benveniste, A., Caillaud, B., et al. (2018). Contracts for system
design. \emph{Foundations and Trends in EDA} 12(2--3).

\bibitem{perrywolf1992} Perry, D.E., Wolf, A.L. (1992). Foundations for the study of software
architecture. \emph{ACM SIGSOFT SEN} 17(4):40--52.

\bibitem{shawgarlan1996} Shaw, M., Garlan, D. (1996). \emph{Software Architecture: Perspectives
on an Emerging Discipline}. Prentice Hall.

\bibitem{boxology1997} Shaw, M., Clements, P. (1997). A field guide to boxology: preliminary
classification of architectural styles for software systems. \emph{Proc.\ COMPSAC 1997}.

\bibitem{lane1990} Lane, T.G. (1990). Studying software architecture through design spaces and
rules. CMU/SEI-90-TR-18; and CMU/SEI-90-TR-22.

\bibitem{shaw2012} Shaw, M. (2012). The role of design spaces. \emph{IEEE Software} 29(1):46--50.

\bibitem{fielding2000} Fielding, R.T. (2000). \emph{Architectural Styles and the Design of
Network-based Software Architectures}. PhD dissertation, UC Irvine.

\bibitem{bass2012} Bass, L., Clements, P., Kazman, R. (2012). \emph{Software Architecture in
Practice}, 3rd ed. Addison-Wesley.

\bibitem{atam2000} Kazman, R., Klein, M., Clements, P. (2000). ATAM: method for architecture
evaluation. CMU/SEI-2000-TR-004.

\bibitem{iso25010} ISO/IEC 25010:2011. Systems and software quality models.

\bibitem{taylor2010} Taylor, R.N., Medvidovic, N., Dashofy, E. (2010). \emph{Software
Architecture: Foundations, Theory, and Practice}. Wiley.

\bibitem{mehta2000} Mehta, N.R., Medvidovic, N., Phadke, S. (2000). Towards a taxonomy of
software connectors. \emph{Proc.\ ICSE 2000}.

\bibitem{kruchten1995} Kruchten, P. (1995). The 4+1 view model of architecture. \emph{IEEE
Software} 12(6):42--50.

\bibitem{evans2003} Evans, E. (2003). \emph{Domain-Driven Design}. Addison-Wesley.

\bibitem{cockburn2005} Cockburn, A. (2005). Hexagonal architecture (ports and adapters).

\bibitem{fowler2014} Fowler, M., Lewis, J. (2014). Microservices. martinfowler.com.

\bibitem{newman2015} Newman, S. (2015). \emph{Building Microservices}. O'Reilly.

\bibitem{jansenbosch2005} Jansen, A., Bosch, J. (2005). Software architecture as a set of
architectural design decisions. \emph{Proc.\ WICSA 2005}.

\bibitem{vangurp2002} van Gurp, J., Bosch, J. (2002). Design erosion: problems and causes.
\emph{Journal of Systems and Software} 61(2):105--119.

\bibitem{kruchten2004} Kruchten, P. (2004). An ontology of architectural design decisions in
software-intensive systems. \emph{Proc.\ 2nd Groningen Workshop on Software Variability}.

\bibitem{maclean1991} MacLean, A., Young, R.M., Bellotti, V., Moran, T.P. (1991). Questions,
options, and criteria: elements of design space analysis. \emph{Human--Computer Interaction}
6(3--4):201--250.

\bibitem{grudin1996} Grudin, J. (1996). Evaluating opportunities for design capture. In Moran \&
Carroll (eds.), \emph{Design Rationale}. Lawrence Erlbaum.

\bibitem{shipman1999} Shipman, F.M., Marshall, C.C. (1999). Formality considered harmful.
\emph{Computer Supported Cooperative Work} 8:333--352.

\bibitem{tang2006} Tang, A., Babar, M.A., Gorton, I., Han, J. (2006). A survey of architecture
design rationale. \emph{Journal of Systems and Software} 79(12):1792--1804.

\bibitem{buchgeher2023} Buchgeher, G., et al. (2023). Using architecture decision records in
open source projects --- an MSR study on GitHub. \emph{IEEE Software}.

\bibitem{bhat2017} Bhat, M., Shumaiev, K., et al. (2017). Automatic extraction of design
decisions from issue management systems. arXiv:1704.04798.

\bibitem{curtis1988} Curtis, B., Krasner, H., Iscoe, N. (1988). A field study of the software
design process for large systems. \emph{Communications of the ACM} 31(11):1268--1287.

\bibitem{zannier2007} Zannier, C., Chiasson, M., Maurer, F. (2007). A model of design decision
making based on empirical results of interviews with software designers. \emph{Information and
Software Technology} 49(6):637--653.

\bibitem{vanvliettang2016} van Vliet, H., Tang, A. (2016). Decision making in software
architecture. \emph{Journal of Systems and Software} 117:638--644.

\bibitem{fowler2003} Fowler, M. (2003). Design --- who needs an architect? \emph{IEEE Software}
20(5):11--13.

\bibitem{bezos2015} Bezos, J. (2015). Amazon shareholder letter (Type 1 / Type 2 decisions).

\bibitem{gardenfors2000} Gärdenfors, P. (2000). \emph{Conceptual Spaces: The Geometry of
Thought}. MIT Press.

\bibitem{shepard1987} Shepard, R.N. (1987). Toward a universal law of generalization for
psychological science. \emph{Science} 237:1317--1323.

\bibitem{tversky1977} Tversky, A. (1977). Features of similarity. \emph{Psychological Review}
84(4):327--352.

\bibitem{hernandez2017} Hernández-Conde, J.V. (2017). A case against convexity in conceptual
spaces. \emph{Synthese} 194:4011--4037.

\bibitem{jager2007} Jäger, G. (2007). The evolution of convex categories. \emph{Linguistics and
Philosophy} 30:551--564.

\bibitem{goelpirolli1992} Goel, V., Pirolli, P. (1992). The structure of design problem spaces.
\emph{Cognitive Science} 16:395--429.

\bibitem{schon1983} Schön, D.A. (1983). \emph{The Reflective Practitioner}. Basic Books.

\bibitem{dorstcross2001} Dorst, K., Cross, N. (2001). Creativity in the design process:
co-evolution of problem--solution. \emph{Design Studies} 22(5):425--437.

\bibitem{gero2004} Gero, J.S., Kannengiesser, U. (2004). The situated
function--behaviour--structure framework. \emph{Design Studies} 25(4):373--391.

\bibitem{rosch1978} Rosch, E. (1978). Principles of categorization. In \emph{Cognition and
Categorization}. Erlbaum.

\bibitem{soloway1984} Soloway, E., Ehrlich, K. (1984). Empirical studies of programming
knowledge. \emph{IEEE TSE} SE-10(5):595--609.

\bibitem{latoza2010} LaToza, T.D., Myers, B.A. (2010). Hard-to-answer questions about code.
\emph{Proc.\ PLATEAU 2010}; and Ko, A.J., DeLine, R., Venolia, G. (2007). Information needs in
collocated software development teams. \emph{Proc.\ ICSE 2007}.

\bibitem{klein1998} Klein, G. (1998). \emph{Sources of Power: How People Make Decisions}. MIT
Press.

\bibitem{dayan1993} Dayan, P. (1993). Improving generalization for temporal difference learning:
the successor representation. \emph{Neural Computation} 5(4):613--624.

\bibitem{stachenfeld2017} Stachenfeld, K.L., Botvinick, M.M., Gershman, S.J. (2017). The
hippocampus as a predictive map. \emph{Nature Neuroscience} 20:1643--1653.

\bibitem{behrens2018} Behrens, T.E.J., et al. (2018). What is a cognitive map? Organizing
knowledge for flexible behavior. \emph{Neuron} 100(2):490--509.

\bibitem{murphy1995} Murphy, G.C., Notkin, D., Sullivan, K. (1995). Software reflexion models:
bridging the gap between source and high-level models. \emph{Proc.\ FSE 1995}.

\bibitem{tzerpos2000} Tzerpos, V., Holt, R.C. (2000). ACDC: an algorithm for
comprehension-driven clustering. \emph{Proc.\ WCRE 2000}.

\bibitem{garcia2011} Garcia, J., Popescu, D., Mattmann, C., Medvidovic, N. (2011). Enhancing
architectural recovery using concerns. \emph{Proc.\ ASE 2011}.

\bibitem{garcia2013} Garcia, J., Ivkovic, I., Medvidovic, N. (2013). A comparative analysis of
software architecture recovery techniques. \emph{Proc.\ ASE 2013}.

\bibitem{lutellier2015} Lutellier, T., et al. (2015). Comparing software architecture recovery
techniques using accurate dependencies. \emph{Proc.\ ICSE 2015}.

\bibitem{louridas2008} Louridas, P., Spinellis, D., Vlachos, V. (2008). Power laws in software.
\emph{ACM TOSEM} 18(1).

\bibitem{tarr1999} Tarr, P., Ossher, H., Harrison, W., Sutton, S.M. (1999). N degrees of
separation: multi-dimensional separation of concerns. \emph{Proc.\ ICSE 1999}.

\bibitem{eaddy2008} Eaddy, M., et al. (2008). Do crosscutting concerns cause defects? \emph{IEEE
TSE} 34(4):497--515.

\bibitem{behnamghader2017} Behnamghader, P., et al. (2017). A large-scale study of architectural
changes in open-source software systems. \emph{Empirical Software Engineering} 22:1146--1193.

\bibitem{azadi2019} Azadi, U., Fontana, F.A., Taibi, D. (2019). Architectural smells detected by
tools: a catalogue proposal. \emph{Proc.\ TechDebt 2019}.

\bibitem{malavolta2013} Malavolta, I., Lago, P., Muccini, H., Pelliccione, P., Tang, A. (2013).
What industry needs from architectural languages. \emph{IEEE TSE} 39(6):869--891.

\bibitem{fitness2017} Ford, N., Parsons, R., Kua, P. (2017). \emph{Building Evolutionary
Architectures}. O'Reilly.

\bibitem{aider2023} Gauthier, P. (2023). Building a better repository map with tree-sitter.
aider.chat.

\bibitem{repograph2024} Ouyang, S., et al. (2024). RepoGraph: enhancing AI software engineering
with repository-level code graphs. arXiv:2410.14684.

\bibitem{contextmatters2026} Gupta, A., Dhar, R., Feitosa, D., Vaidhyanathan, K. (2026). Context
matters: evaluating context strategies for automated ADR generation using LLMs. arXiv:2604.03826.

\bibitem{gamma1994} Gamma, E., Helm, R., Johnson, R., Vlissides, J. (1994). \emph{Design
Patterns}. Addison-Wesley.

\bibitem{tanenbaum1992} Tanenbaum, A., Torvalds, L. (1992). LINUX is obsolete
(comp.os.minix debate). In \emph{Open Sources} (O'Reilly, 1999), appendix.

\bibitem{chromium2009} Reis, C., Gribble, S.D. (2009). Isolating web programs in modern browser
architectures. \emph{Proc.\ EuroSys 2009}.

\bibitem{reis2019} Reis, C., Moshchuk, A., Oskov, N. (2019). Site isolation: process separation
for web sites within the browser. \emph{Proc.\ USENIX Security 2019}.

\bibitem{weiss2024} Ahmeti, B., Linder, M., Groner, R., Wohlrab, R. (2024). Architecture decision
records in practice: an action research study. \emph{Proc.\ ECSA 2024}, LNCS 14889:333--349.

\bibitem{yourdon1979} Yourdon, E., Constantine, L.L. (1979). \emph{Structured Design}.
Prentice Hall; and Myers, G.J. (1975). \emph{Reliable Software through Composite Design}.
Petrocelli/Charter.

\bibitem{sqlitearch} Hipp, D.R. Architecture of SQLite. sqlite.org/arch.html; and Gaffney, K.P.,
Prammer, M., Brasfield, L., Hipp, D.R., Kennedy, D., Patel, J.M. (2022). SQLite: past, present,
and future. \emph{Proceedings of the VLDB Endowment} 15(12):3535--3547.

\bibitem{fiber2016} Clark, A. (acdlite) (2016). React fiber architecture.
github.com/acdlite/react-fiber-architecture.

\bibitem{poettering2010} Poettering, L. (2010). Rethinking PID 1. 0pointer.de.

\bibitem{primevideo2023} Kolny, M. (2023). Scaling up the Prime Video audio/video monitoring
service and reducing costs by 90\%. Prime Video Tech Blog.

\bibitem{segment2018} Noonan, A. (2018). Goodbye microservices. Segment Engineering Blog.

\bibitem{shopify2019} Westeinde, K. (2019). Deconstructing the monolith. Shopify Engineering.

\bibitem{uberdoma2020} Gluck, A. (2020). Introducing domain-oriented microservice architecture.
Uber Engineering.

\bibitem{istio2020} Istio team (2020). Introducing istiod: simplifying the control plane.
istio.io.

\bibitem{luo2021} Luo, S., et al. (2021). Characterizing microservice dependency and
performance: Alibaba trace analysis. \emph{Proc.\ SoCC 2021}.

\bibitem{aosanginx} Alexeev, A. nginx. In \emph{The Architecture of Open Source Applications},
vol.~II. aosabook.org.

\bibitem{aosagit} Potter, S. Git. In \emph{AOSA}, vol.~II. aosabook.org.

\bibitem{aosaeclipse} Moir, K. Eclipse. In \emph{AOSA}, vol.~I. aosabook.org.

\bibitem{aosazeromq} Sústrik, M. ZeroMQ. In \emph{AOSA}, vol.~II. aosabook.org.

\bibitem{hunt2010} Hunt, P., Konar, M., Junqueira, F.P., Reed, B. (2010). ZooKeeper: wait-free
coordination for internet-scale systems. \emph{Proc.\ USENIX ATC 2010}.

\bibitem{c4model} Brown, S. The C4 model for visualising software architecture. c4model.com.

\bibitem{madr2018} Kopp, O., Armbruster, A., Zimmermann, O. (2018). Markdown architectural
decision records: format and tool support. \emph{Proc.\ ZEUS 2018}, CEUR-WS vol.~2072.

\bibitem{postgres2023} Linnakangas, H. (2023). Let's make PostgreSQL multi-threaded.
pgsql-hackers mailing list, postgresql.org.

\end{thebibliography}
\end{document}
